% !TeX root = ../complete.tex
\section{کاهش \lr{Switching Power}}
\label{sec:power-dynamic-reduction}

از رابطه \eqref{eq:power-switching} چهار کمیت برای کاهش توان به دست
می‌آیند. توان با \(\alpha\)، \(C_L\) و \(f\) متناسب است و با
مربع \(V_{DD}\) تغییر می‌کند. شرط این مقایسه آن است که هنگام تغییر
هر کمیت، بقیه ثابت باشند. در طراحی واقعی، تغییر ظرفیت یا ولتاژ می‌تواند
تأخیر مسیر را نیز تغییر دهد. پس بعد از برآورد توان، باید معتبر بودن
فرض‌ها و برقراری محدودیت زمانی را بررسی کرد
\cite[بخش \secnum{۵.۲}، صفحات ۱۸۵ تا ۱۹۳]{weste}
\cite[اسلایدهای ۱۳ و ۱۷ تا ۱۹]{powerslides}.

\subsection{توقف ساعت یک بلوک بی‌کار}
\label{sec:power-clock-gating}

\lr{Register} یک یا چند بیت را نگه می‌دارد. در این مثال، در هر لبه
صعودی ساعت، داده ورودی را می‌گیرد و تا لبه بعد نگه می‌دارد. اگر
\lr{Register} ورودی یک بلوک در چند چرخه به‌روزرسانی لازم نداشته باشد،
با \lr{Clock Gating} می‌توان لبه‌های ساعتِ همان بلوک را در آن مدت
متوقف کرد. در این حالت هم گذارهای شاخه ساعت و هم گذارهایی که از
تغییر خروجی این \lr{Register} ناشی می‌شوند، کاهش می‌یابند. اگر
ورودی‌های دیگری مستقل از آن \lr{Register} تغییر کنند، ممکن است بخشی
از منطق بلوک همچنان گذار داشته باشد
\cite[بخش \secnum{۵.۲.۱.۱}، صفحه ۱۸۶]{weste}
\cite[اسلاید ۱۷]{powerslides}.

برای فهم مدار، ساعت اصلی را \(\mathrm{clk}\)، سیگنال اجازه عبور را
\(e\) و ساعت بلوک را \(\mathrm{gclk}\) بنامید. در مدل منطقی،
\(\mathrm{gclk}=\mathrm{clk}\land e\) است. وقتی \(e=1\)، ساعت
بلوک از ساعت اصلی پیروی می‌کند و وقتی \(e=0\)، ساعت بلوک صفر
می‌ماند. اگر \(e\) در میانه سطح بالای ساعت تغییر کند، این گیت
می‌تواند پالس را کوتاه کند یا لبه‌ای ناخواسته به \lr{Register}
بفرستد. راه شکل \secnum{۵.۷} کتاب استفاده از \lr{Latch} است،
یعنی عنصر حافظه‌ای که وقتی ساعت صفر است مقدار اجازه عبور را می‌پذیرد
و در تمام زمان یک‌بودن ساعت آن را ثابت نگه می‌دارد
\cite[شکل \secnum{۵.۷}، صفحه ۱۸۶]{weste}.

اگر یک شاخه در ۱۰ دوره کامل ساعت اصلی تنها چهار دوره پالس داشته
باشد، \(\alpha\) آن شاخه نسبت به ساعت اصلی \(4/10=0.4\) است.
با ثابت بودن ولتاژ و ظرفیت شاخه، توان سوئیچینگ این شاخه ۴۰ درصد
حالت بدون \lr{Clock Gating} می‌شود. اندازه‌گیری صرفه‌جویی کل بلوک
به فعالیت داده‌ها، مصرف مدار کنترل و اجزایی که همچنان ساعت می‌گیرند
نیز وابسته است. تغذیه مدار در این روش وصل می‌ماند.

\subsection{کاهش ظرفیت و کنترل تأخیر}
\label{sec:power-capacitance-reduction}

کوتاه کردن سیمِ گره‌هایی که زیاد سوئیچ می‌کنند، سهم ظرفیت سیم را در
رابطه \eqref{eq:power-node-sum} کم می‌کند. کوچک کردن یک گیت، ظرفیت
ورودی آن را برای مرحله پیشین و ظرفیت \lr{Diffusion} خود گیت را نیز
کم می‌کند. در عوض، گیت کوچک‌تر معمولاً همان بار بعدی را کندتر شارژ
می‌کند. در یک مسیر، اثر هر تغییر را روی ظرفیت، تأخیر مراحل پیش و پس
از گیت و فعالیت گره‌ها بررسی می‌کنیم. مبنای بررسی تأخیر همان
\NotesTopic{sec:delay-path-model}{مدل مسیر در \lr{Logical Effort}}
است \cite[بخش \secnum{۵.۲.۲}، صفحات ۱۸۸ تا ۱۹۰]{weste}
\cite[اسلاید ۱۸]{powerslides}.

برای یک سیم بلند، یکی از گزینه‌ها این است که یک
\lr{Inverter} یا \lr{Buffer} در ابتدای سیم قرار دهیم تا
ظرفیت بزرگ سیم مستقیماً به خروجی گیت منطقی تحمیل نشود. در این
آرایش، گیت منطقی ورودی مرحله جدید را راه‌اندازی می‌کند و
مرحله جدید سیم را شارژ و تخلیه می‌کند. با افزایش اندازه
\lr{Buffer}، بار ورودی آن نیز بیشتر می‌شود. برای تصمیم‌گیری،
ظرفیت ورودی و ظرفیت‌های درونی مرحله اضافه، قطبیت خروجی و تأخیر
کل مسیر را در کنار ظرفیت سیم در نظر می‌گیریم
\cite[بخش \secnum{۵.۲.۲}، صفحه ۱۸۸]{weste}
\cite[اسلاید ۱۸]{powerslides}.

برای نمونه، دو \lr{Inverter} پیاپی با ظرفیت‌های ورودی ۱ و ۴ واحد
و بار نهایی ۱۶ واحد در نظر بگیرید. واحد ظرفیت را ظرفیت ورودی
\lr{Inverter} مرجع و \(\tau\) را واحد تأخیر فصل قبل می‌گیریم.
در مدل عادی‌شده فصل تأخیر، تأخیر هر \lr{Inverter} شامل یک واحد
سهم درونی و نسبت ظرفیت بار به ظرفیت ورودی آن است. با چشم‌پوشی از
ظرفیت سیم و شاخه‌های فرعی، فقط دو گیت داریم. اگر ظرفیت ورودی گیت
دوم را به ۲ واحد برسانیم، تأخیر کل به‌ترتیب قبل و بعد از تغییر چنین است.
\[
  \begin{aligned}
    D_{\mathrm{before}} &= \left(1+\frac{4}{1}\right)
                         +\left(1+\frac{16}{4}\right)=10,\\
    D_{\mathrm{after}}  &= \left(1+\frac{2}{1}\right)
                         +\left(1+\frac{16}{2}\right)=12.
  \end{aligned}
\]
برای تبدیل این اعداد به تأخیر، هرکدام را در \(\tau\) ضرب می‌کنیم.
تأخیر نرمال‌شده گیت اول از ۵ به ۳ کم شده، اما
تأخیر گیت دوم از ۵ به ۹ رسیده است. گیت کوچک‌تر ظرفیت ورودی کم‌تری
دارد و با فرض فعالیت و دامنه یکسان، انرژی شارژ مربوط به همین بخش
از بار گره میانی نصف می‌شود. اگر مهلت مسیر کم‌تر از \(12\tau\)
باشد، این انتخاب با وجود کاهش انرژیِ بار میانی پذیرفتنی نیست.

\subsection{تنظیم ولتاژ و فرکانس}
\label{sec:power-voltage-frequency}

برای مقایسه دو حالت کاریِ یک گره، \(V_1,f_1\) و \(V_2,f_2\)
را ولتاژ و فرکانس حالت‌های اول و دوم بگیرید. اگر ظرفیت و فعالیت
به‌ازای هر دوره ثابت باشند و توان سوئیچینگ حالت اول مثبت باشد، با
تقسیم رابطه \eqref{eq:power-switching} برای دو حالت، عامل‌های مشترک
ساده می‌شوند و نسبت زیر به دست می‌آید.
\begin{equation}
  \frac{P_{\mathrm{sw},2}}{P_{\mathrm{sw},1}}
    =\left(\frac{V_2}{V_1}\right)^2\frac{f_2}{f_1}.
  \label{eq:power-voltage-frequency-ratio}
\end{equation}
مثلاً اگر ولتاژ به ۸۰ درصد و فرکانس به نصف مقدار اولیه برسد، نسبت
توان \(0.8^2\times0.5=0.32\) می‌شود. برای یک کار با تعداد دوره
ثابت، نصف شدن فرکانس زمان انجام کار را دو برابر می‌کند. در این مثال
نسبت انرژی سوئیچینگِ آن کار \(0.32\times2=0.64\) است. اگر ولتاژ
ثابت بماند، کاهش فرکانس به‌تنهایی توان سوئیچینگ را کم می‌کند، ولی
با همین فرض‌ها انرژی سوئیچینگِ کار ثابت می‌ماند. نتیجه به مؤلفه
سوئیچینگ محدود است و انرژی نشتی در زمان اجرای طولانی‌تر می‌تواند
تغییر کند \cite[بخش‌های \secnum{۵.۲.۳} و \secnum{۵.۲.۴}،
صفحات ۱۹۰ تا ۱۹۳]{weste}.

ولتاژ کم‌تر معمولاً جریان راه‌اندازی را کم و گذار را کندتر می‌کند.
ساعت را باید چنان انتخاب کرد که مدار با این تأخیر هنوز کار را در
زمان مقرر انجام دهد. در \lr{Dynamic Voltage and Frequency Scaling}
یا \lr{DVFS}، این دو کمیت برحسب نیاز عملکردیِ بار کاری تنظیم
می‌شوند. در طراحی با \lr{Voltage Domain}، هر بلوک می‌تواند تغذیه
مختص خود را داشته باشد. وقتی سیگنالی از تغذیه کم‌تر به بیش‌تر می‌رود،
سطح یکِ فرستنده ممکن است برای خاموش‌کردن کامل \lr{pMOS} گیرنده
کافی نباشد. \lr{Level Converter} دامنه سیگنال را برای بلوک گیرنده
بازسازی می‌کند. زمان و انرژی خود مبدل نیز در مقایسه معماری‌ها سهم
دارند \cite[بخش \secnum{۵.۲.۳}، صفحات ۱۹۰ تا ۱۹۲]{weste}
\cite[اسلاید ۱۹]{powerslides}.

در \lr{Resonant Clock}، کتاب راه دیگری را برای بار دوره‌ای معرفی
می‌کند. انرژی میان خازن و سلف، یعنی عنصری که انرژی را در میدان
مغناطیسی ذخیره می‌کند، مبادله می‌شود. راه‌انداز در هر دوره انرژی
ازدست‌رفته در مقاومت‌ها را جبران می‌کند. این تغییر مسیر انرژی، فرض
مدل معمول شارژ از منبع و تخلیه کامل به زمین را عوض می‌کند. در برابر
امکان کاهش مصرف، مدار به اجزای بیشتر و فرکانس کاری متناسب با
شبکه تشدید نیاز دارد \cite[بخش \secnum{۵.۲.۶}، صفحه ۱۹۳]{weste}.
